\documentclass[11pt]{article}

\usepackage[T1]{fontenc}
\usepackage{lmodern}
\usepackage[margin=1in]{geometry}
\usepackage{amsmath,amssymb,amsthm}
\usepackage{booktabs}
\usepackage[hidelinks]{hyperref}

\newtheorem{theorem}{Theorem}
\newtheorem{lemma}[theorem]{Lemma}
\newtheorem{proposition}[theorem]{Proposition}
\newcommand{\Aut}{\operatorname{Aut}}

\title{The Minimum 2-Circuit Decomposition of an
18-Triangle Torus}
\author{Xiaojian Zeng\\
PBC School of Finance, Tsinghua University, Beijing 100190, China\\
\texttt{xiaojianzeng@foxmail.com}}
\date{September 2026}

\hypersetup{
  pdftitle={The Minimum 2-Circuit Decomposition of an 18-Triangle Torus},
  pdfauthor={Xiaojian Zeng}
}

\begin{document}
\maketitle

\begin{abstract}
For the nine-vertex triangulated torus displayed by Harlander and
Quatrone, we prove that the minimum number of 2-circuits in a facet-exact
decomposition is three.  The proof reduces every possible circuit to an orbit
of a forced transition permutation on the 108 ordered facet states.  The
transition orbits give exactly nine circuits up to rotation and reversal, and
their disjointness graph gives exactly three labeled minimum decompositions.
An exhaustive calculation of the automorphism group shows that these three
decompositions form one orbit.  The finite enumeration is supplied with a
standard-library verifier that independently reconstructs all incidence,
transition, exact-cover, and automorphism calculations.
\end{abstract}

\section{The finite problem and result}

Let $V=\{0,1,\ldots,8\}$, and let $S$ be the abstract simplicial complex
whose facets, in the order used below, are
\begin{equation}\label{eq:facets}
\begin{split}
F(S)=\{&105,053,537,374,741,410,\\
       &216,165,658,587,872,721,\\
       &023,236,364,648,480,802\}.
\end{split}
\end{equation}
Each three-digit string denotes an unordered triple.  This is the torus in
Figure~1 of Harlander and Quatrone~\cite{HarlanderQuatrone2026}.  That figure
gives the three 2-circuits
\begin{equation}\label{eq:source-cover}
105374,\qquad 216587,\qquad 023648
\end{equation}
and leaves the nonexistence of a single Euler 2-circuit to the reader.  It does
not give an exclusion of a two-circuit decomposition or the census below.

A \emph{2-circuit} in the present finite complex is a cyclic word
$W=(v_i)_{i\in\mathbb Z/n\mathbb Z}$ such that every cyclic window
$\{v_i,v_{i+1},v_{i+2}\}$ is a facet and these $n$ facets are pairwise
distinct.  A \emph{2-circuit decomposition} is an unordered collection of
such words whose facet sets partition $F(S)$.  Write $c(S)$ for the minimum
number of circuits in such a decomposition.

This consecutive-window viewpoint is part of the theory of universal
cycles, developed by Chung, Diaconis, and Graham for families of
combinatorial objects~\cite{ChungDiaconisGraham1992}.  In hypergraph
language, Glock, Joos, K\"uhn, and Osthus proved the existence of tight Euler
tours in quasirandom uniform hypergraphs under the natural divisibility
condition; their complete-hypergraph case settles the corresponding
universal-cycle conjecture for fixed-size subsets~\cite{GlockEtAl2020}.
The facet hypergraph in \eqref{eq:facets} is instead one fixed sparse surface
triangulation.  Those general existence results neither determine its minimum
number of circuits nor give its exact decomposition census.

Two words representing the same circuit are identified under cyclic rotation
and reversal.  Decompositions are also identified under permutation of their
circuits.  For the final quotient, $\Aut(S)$ acts diagonally by relabeling
vertices.

\begin{theorem}\label{thm:main}
For the complex in \eqref{eq:facets},
\[
c(S)=3.
\]
The words in \eqref{eq:source-cover} form a minimum decomposition.  Up to
rotation and reversal of the individual circuits and permutation of the
circuits, there are exactly three labeled minimum decompositions.  The group
$\Aut(S)$ has order $108$ and is transitive on the three decompositions;
hence there is exactly one minimum-decomposition orbit after quotienting by
$\Aut(S)$.
\end{theorem}

The theorem concerns only the fixed complex \eqref{eq:facets}; it makes no
claim about arbitrary triangulated tori.

\section{A complete transition reduction}

Let
\[
\Omega=\{(a,b,c):\{a,b,c\}\in F(S)\}.
\]
Thus every facet contributes its six orderings and $|\Omega|=108$.  Directly
from \eqref{eq:facets}, $S$ has $27$ edges and every edge lies in exactly two
facets.  Define $P:\Omega\to\Omega$ by
\begin{equation}\label{eq:transition}
P(a,b,c)=(b,c,d),
\end{equation}
where $\{b,c,d\}$ is the other facet containing the edge $\{b,c\}$.
This is a permutation.  Indeed, its inverse sends $(x,y,z)$ to $(w,x,y)$,
where $\{w,x,y\}$ is the other facet at $\{x,y\}$.

\begin{lemma}[Orbit completeness]\label{lem:complete}
The 2-circuits in $S$ are precisely the $P$-orbits whose projections to
unordered facets are injective.
\end{lemma}

\begin{proof}
Let $W=(v_i)$ be a 2-circuit and put
$s_i=(v_i,v_{i+1},v_{i+2})$.  The unordered facets represented by $s_i$ and
$s_{i+1}$ are distinct and share $\{v_{i+1},v_{i+2}\}$.  Since exactly two
facets contain this edge, \eqref{eq:transition} forces
$s_{i+1}=P(s_i)$.  Cyclic closure gives $P^n(s_0)=s_0$.  The word cannot stop
before the state orbit closes, and it cannot traverse an orbit more than once
without repeating a facet.  Hence its states form one complete $P$-orbit.
Conversely, the overlapping entries along a $P$-orbit form a cyclic word; if
its facet projection is injective, that word is a 2-circuit.
\end{proof}

This lemma is the step that turns the computation into a complete search:
there is no unbounded word length or unexplored prefix after the 108 states
have been exhausted.

\section{Exact circuit and cover enumeration}

Index the facets in \eqref{eq:facets} from $0$ through $17$.  Iterating $P$
gives 18 directed orbits, all of length six and all with injective facet
projection.  Reversal pairs them into the nine circuits in
Table~\ref{tab:circuits}.  The table is also a compact checkable certificate:
for each displayed cyclic word, its six length-three windows and the six
windows of its reversal are two $P$-orbits.  Across the nine rows these are
108 distinct ordered states, so the list exhausts $\Omega$.

\begin{table}[ht]
\centering
\caption{All circuits modulo cyclic rotation and reversal.}
\label{tab:circuits}
\begin{tabular}{@{}cll@{}}
\toprule
Index & Cyclic word & Facet indices \\
\midrule
$C_1$ & $014735$ & $0,1,2,3,4,5$ \\
$C_2$ & $015684$ & $0,5,7,8,15,16$ \\
$C_3$ & $023648$ & $12,13,14,15,16,17$ \\
$C_4$ & $028753$ & $1,2,9,10,12,17$ \\
$C_5$ & $032615$ & $0,1,6,7,12,13$ \\
$C_6$ & $041728$ & $4,5,10,11,16,17$ \\
$C_7$ & $126347$ & $3,4,6,11,13,14$ \\
$C_8$ & $127856$ & $6,7,8,9,10,11$ \\
$C_9$ & $346857$ & $2,3,8,9,14,15$ \\
\bottomrule
\end{tabular}
\end{table}

The disjointness graph of the rows in Table~\ref{tab:circuits} is the disjoint
union of three triangles,
\[
\{C_1,C_3,C_8\},\qquad
\{C_2,C_4,C_7\},\qquad
\{C_5,C_6,C_9\}.
\]
Each triangle partitions all 18 facets, so the three labeled decompositions
are
\begin{align*}
&\{014735,023648,127856\},\\
&\{015684,028753,126347\},\\
&\{032615,041728,346857\}.
\end{align*}
Every possible circuit has six facets.  Consequently, one or two circuits can
cover at most 12 of the 18 facets, while the first displayed decomposition is
equivalent under rotation and reversal to \eqref{eq:source-cover}.  This proves
$c(S)=3$ and proves the labeled census.

For the automorphism quotient, every one of the $9!$ permutations of $V$ was
tested against the facet set \eqref{eq:facets}; exactly 108 preserve it.  Their
induced action on the three covers realizes each element of $S_3$ exactly 18
times.  In particular, the automorphisms written as images of
$012345678$ by
\[
012854763 \quad\text{and}\quad 034125768
\]
send the first cover to the second and third, respectively.  Thus the action
is transitive.  Equivalently, the Burnside fixed-point sum is 108 and the
number of orbits is $108/108=1$.

\section{Independent finite verification}

The ancillary file \texttt{circuit\_certificate.json} records the edge
incidence table, the transition orbits, the nine unoriented circuits, the
three exact covers, all 108 automorphisms, and their action on the covers.
The program \texttt{replay\_certificate.py} does not import the program that
generated this file.  Starting only from the 18 facets, it independently
rebuilds the edge incidences and transition permutation, walks from every
ordered facet, tests all $2^9=512$ circuit subcollections with exact facet
multiplicities, exhausts all $9!$ vertex permutations, and compares the
result with the certificate.  It uses only integer and finite-set operations
from the Python standard library.

A related certificate architecture---exact finite predicates, exhaustive
symmetry calculations, and a separate verifier that reconstructs the
theorem-bearing objects---is documented by Zeng, Liu, Ratnavelu, and
Huat~\cite{ZengEtAl2026}.  Here the human-readable orbit-completeness lemma is
essential: successful execution alone is not used as a substitute for the
mathematical reduction.

To replay the present calculation from the project root, run
\begin{verbatim}
python3 evidence/replay_certificate.py
\end{verbatim}
The expected terminal line is \texttt{REPLAY OK}.  No floating-point
arithmetic, random seed, or numerical tolerance occurs in the proof.

\section*{Computational and AI assistance disclosure}

OpenAI Codex assisted with literature search, finite search design, draft
code, adversarial audit organization, and initial manuscript preparation.
The mathematical claims rest on the explicit reduction above and the exact
ancillary certificate and verifier, not on language-model output.  No proof
assistant was used.

\bibliographystyle{plain}
\bibliography{references}

\end{document}
